\documentclass[aspectratio=169, 10pt]{beamer}

\usetheme{Madrid}
\usecolortheme{default}
\definecolor{ucbblue}{RGB}{0, 51, 102}

\setbeamercolor{structure}{fg=ucbblue}
\setbeamercolor*{palette primary}{fg=white,bg=ucbblue}
\setbeamercolor*{palette secondary}{fg=white,bg=ucbblue!80}
\setbeamercolor*{palette tertiary}{fg=white,bg=ucbblue!90}
\setbeamercolor{title}{fg=white, bg=ucbblue}
\setbeamercolor{frametitle}{fg=white, bg=ucbblue}
\setbeamercolor{item}{fg=ucbblue}

\usepackage[T1]{fontenc}
\usepackage[utf8]{inputenc}
\usepackage{tgpagella}
\usefonttheme{serif}
\usepackage[spanish, es-noshorthands]{babel}
\usepackage{amsmath, amssymb, bm}
\usepackage{booktabs}
\usepackage[most]{tcolorbox}
\usepackage{tikz}
\usepackage{pgfplots}
\pgfplotsset{compat=1.18}
\usepackage[american]{circuitikz}

\title[Consolidación Guiada RC]{Unidad EC3: Consolidación Guiada de Transitorios RC}
\subtitle{Revisión de Ejercicios y Verificación Analítica de Autonomía}
\author[B. Quiroga Turdera]{M.Sc. Ing. Bernardo Quiroga Turdera}
\institute[UCB Tarija]{Universidad Católica Boliviana ``San Pablo'' --- Sede Tarija}
\date{Semana 11}

\begin{document}

\begin{frame}
    \titlepage
\end{frame}

\begin{frame}{Revisión Diagnóstica: Evidencia de Autonomía}
    \textbf{Auditoría de Ejercicios Asignados}
    \begin{itemize}
        \item El objetivo no es tener la "respuesta correcta" copiada, sino demostrar el \textbf{flujo lógico}.
        \item Un error numérico al final de la fórmula $\neq$ un fallo en la física del circuito.
    \end{itemize}
    \vspace{0.4cm}
    \begin{tcolorbox}[colback=white, colframe=ucbblue, title=Preguntas Centrales de Revisión]
        \begin{enumerate}
            \item ¿Qué valor tiene $v_C(0^-)$?
            \item ¿Por qué afirmamos que $v_C(0^+) = v_C(0^-)$?
            \item ¿Cuál es la topología del circuito para $t>0$?
            \item ¿Qué ley física introduce el término derivativo $\frac{dv_C}{dt}$?
        \end{enumerate}
    \end{tcolorbox}
\end{frame}

\begin{frame}{Paso 1: La Condición Inicial ($v_C(0^-)$ y $v_C(0^+)$)}
    \begin{columns}[T]
        \column{0.5\textwidth}
        \textbf{Inercia Eléctrica}
        \begin{itemize}
            \item Bajo condiciones ordinarias, el voltaje en un capacitor ideal \textbf{no salta instantáneamente}.
            \item Requeriría un pulso infinito de corriente: $i_C = C \frac{dv_C}{dt}$.
        \end{itemize}
        \begin{alertblock}{Principio de Continuidad}
            $$v_C(0^+) = v_C(0^-)$$
        \end{alertblock}

        \column{0.5\textwidth}
        \textbf{Metodología:}
        \begin{enumerate}
            \item Analizar el circuito en $t < 0$.
            \item Determinar la tensión de equilibrio ($v_C(0^-)$).
            \item \textbf{Heredar} ese valor exacto para el instante $t = 0^+$, sin importar cómo cambie el circuito después del interruptor.
        \end{enumerate}
    \end{columns}
\end{frame}

\begin{frame}{Paso 2: La Condición Final ($v_C(\infty)$)}
    \begin{columns}[T]
        \column{0.5\textwidth}
        \textbf{Estado Estacionario DC}
        \begin{itemize}
            \item Analizamos la red \textbf{después de la conmutación} ($t > 0$) cuando ha pasado mucho tiempo.
            \item La carga termina, la tensión deja de variar: $\frac{dv_C}{dt} \to 0$.
            \item Por lo tanto, $i_C \to 0$. El capacitor actúa como una \textbf{rama abierta}.
        \end{itemize}
        
        \column{0.5\textwidth}
        \textbf{Cálculo Práctico:}
        \begin{center}
        \begin{circuitikz}[scale=0.8, transform shape]
            \draw (0,0) to[V, l={$12\,\text{V}$}] (0,3) to[R, l={$2\,\text{k}\Omega$}] (3,3);
            \draw (3,3) to[short, *-o] (4.5,3) node[above]{$a$};
            \draw (3,3) to[R, l={$4\,\text{k}\Omega$}] (3,0);
            \draw (4.5,3) to[open, v={$v_C(\infty)$}] (4.5,0);
            \draw (0,0) -- (4.5,0) node[ground]{};
        \end{circuitikz}
        \end{center}
        \vspace{-0.3cm}
        Divisor de Tensión:
        $$v_C(\infty) = 12\,\text{V} \cdot \left( \frac{4\,\text{k}\Omega}{2\,\text{k}\Omega + 4\,\text{k}\Omega} \right) = \mathbf{8\,\text{V}}$$
    \end{columns}
\end{frame}

\begin{frame}{Paso 3: Resistencia Equivalente vista por el Capacitor ($R_{eq}$)}
    \textbf{La Perspectiva del Almacenador}
    \begin{itemize}
        \item La rapidez con la que reacciona el circuito depende de la resistencia que ve el capacitor al mirar hacia la red ($t > 0$).
        \item \textbf{Regla:} Retirar el capacitor y \textbf{desactivar} las fuentes independientes (cortocircuito para fuentes de voltaje, circuito abierto para fuentes de corriente).
    \end{itemize}
    
    \begin{center}
    \begin{circuitikz}[scale=0.8, transform shape]
        \draw (0,3) to[short, -*] (3,3);
        \draw (0,0) to[short] (0,3); % Cortocircuito de la fuente de 12V
        \draw (3,3) to[R, l_={$2\,\text{k}\Omega$}] (0,3);
        \draw (3,3) to[R, l_={$4\,\text{k}\Omega$}] (3,0);
        \draw (3,3) -- (5,3) node[right]{$+$};
        \draw (3,0) -- (5,0) node[right]{$-$};
        \draw (0,0) -- (3,0) node[ground]{};
        \node at (5.5, 1.5) {$R_{eq} = 2\,\text{k}\Omega \parallel 4\,\text{k}\Omega = 1.333\,\text{k}\Omega$};
    \end{circuitikz}
    \end{center}
\end{frame}

\begin{frame}{Paso 4: Constante de Tiempo ($\tau$) y Forma Exponencial}
    \begin{columns}[T]
        \column{0.4\textwidth}
        \textbf{Constante de Tiempo:}
        \begin{block}{}
            $$\tau = R_{eq} \cdot C$$
        \end{block}
        \begin{itemize}
            \item Unidades: Segundos [s].
            \item Mayor $\tau \implies$ Respuesta más lenta.
        \end{itemize}
        
        \column{0.6\textwidth}
        \textbf{Ecuación Universal (1er Orden):}
        \begin{tcolorbox}[colback=white, colframe=ucbblue, boxrule=0.5pt, arc=1mm]
            $$v_C(t) = v_C(\infty) + \left[ v_C(0^+) - v_C(\infty) \right] e^{-t/\tau}$$
        \end{tcolorbox}
        \begin{itemize}
            \item $v_C(\infty)$: Destino (Asíntota).
            \item $[v_C(0^+) - v_C(\infty)]$: Desviación inicial.
            \item \textbf{Verificación:} Si $t=0 \implies e^0=1 \implies v_C(0^+)$. Si $t \to \infty \implies e^{-\infty}=0 \implies v_C(\infty)$.
        \end{itemize}
    \end{columns}
\end{frame}

\begin{frame}{Ejemplo Guiado de Transferencia Analítica}
    \begin{columns}[T]
        \column{0.4\textwidth}
        En $t=0$, el circuito entra en acción. Se sabe que $v_C(0^-) = 0\,\text{V}$.
        \vspace{0.2cm}
        \begin{center}
        \begin{circuitikz}[scale=0.8, transform shape]
            \draw (0,3) to[V, l_={$9\,\text{V}$}] (0,0);
            \draw (0,3) to[closing switch, l={$t=0$}] (2,3) to[R, l={$3\,\text{k}\Omega$}] (4,3);
            \draw (4,3) to[short, *-o] (5.5,3) node[above]{$a$};
            \draw (4,3) to[R, l_={$6\,\text{k}\Omega$}] (4,0);
            \draw (5.5,3) to[C, l={$100\,\mu\text{F}$}, v={$v_C$}] (5.5,0);
            \draw (0,0) -- (5.5,0) node[ground]{};
        \end{circuitikz}
        \end{center}
        
        \column{0.6\textwidth}
        \begin{enumerate}
            \item \textbf{Inicial:} $v_C(0^+) = v_C(0^-) = \mathbf{0\,\text{V}}$.
            \item \textbf{Final:} Divisor $9\,\text{V} \cdot \left( \frac{6\text{k}\Omega}{3\text{k}\Omega + 6\text{k}\Omega} \right) = \mathbf{6\,\text{V}}$.
            \item \textbf{$\mathbf{R_{eq}}$:} $3\text{k}\Omega \parallel 6\text{k}\Omega = \mathbf{2\,\text{k}\Omega}$.
            \item \textbf{$\tau$:} $(2000\,\Omega)(100\mu\text{F}) = \mathbf{0.2\,\text{s}}$.
            \item \textbf{Ecuación:} $v_C(t) = 6 + [0 - 6]e^{-t/0.2}$ \\
            $\mathbf{v_C(t) = 6(1 - e^{-5t})\,\text{V}}$
        \end{enumerate}
    \end{columns}
\end{frame}

\begin{frame}{Puente hacia la Simulación (\texttt{.tran}) - Próxima Sesión}
    La simulación en SPICE no genera mágicamente resultados; confirma nuestra predicción analítica de cómo el circuito administra la inercia temporal.
    
    \vspace{0.4cm}
    \renewcommand{\arraystretch}{1.5}
    \begin{tabular}{p{4.5cm} p{8cm}}
    \toprule
    \textbf{Predicción Analítica} & \textbf{Validación Futura Computacional} \\ 
    \midrule
    $v_C(0^+)$ (Condición inicial) & Valor inicial exacto del trazo de la señal en $t=0$. \\ 
    $v_C(\infty)$ (Condición final) & Asíntota u offset DC al final de la simulación. \\ 
    Constante de tiempo $\tau$ & La rapidez visible. Asentamiento de la curva $\sim 5\tau$. \\ 
    Tendencia de la exponencial & Sube o baja suavemente, sin brincos instantáneos de voltaje. \\ 
    \bottomrule
    \end{tabular}
\end{frame}

\end{document}
